% figures/fig1_discretizacion.tex
\begin{tikzpicture}[
    >=Stealth,
    % Cajas de la columna izquierda (física instrumental)
    block/.style={
        rectangle,
        draw=black!70,
        fill=black!2,
        rounded corners=4pt,
        line width=0.8pt,
        text width=5.8cm,
        align=center,
        inner sep=8pt,
        minimum height=1.35cm,
        font=\small\sffamily
    },
    % Fórmulas matemáticas de la columna derecha
    state/.style={
        font=\small,
        align=center,
        inner sep=6pt
    },
    % Flechas principales
    arrow/.style={
        ->,
        line width=1.05pt,
        draw=black!75,
        shorten >= 2pt,
        shorten <= 2pt
    },
    % Líneas de unión sin punta
    line/.style={
        line width=1.05pt,
        draw=black!75
    },
    % Operadores circulares
    oplabel/.style={
        circle,
        draw=black!85,
        fill=white,
        line width=0.9pt,
        inner sep=3.5pt,
        font=\normalsize\bfseries
    },
    % Cajas terminales de las ramas (Parte 2)
    branchbox/.style={
        rectangle,
        draw=UdeAGreen!85!black,
        fill=UdeAGreen!6,
        rounded corners=4pt,
        line width=0.9pt,
        text width=5.4cm,
        align=center,
        inner sep=7pt,
        font=\small\sffamily
    }
]

    % =========================================================================
    % Nivel 0: Campo Continuo (y = 0.0)
    % =========================================================================
    \node[block] (b0) at (0, 0) {
        \textbf{Campos acústicos continuos}\\
        \footnotesize Presión sonora en receptores duales
    };
    \node[state] (s0) at (7.2, 0) {
        $\mathbf{p}(t) = \begin{pmatrix} p_1(t) \\ p_2(t) \end{pmatrix} \in L^2(\mathbb{R}) \times L^2(\mathbb{R})$
    };

    % Operador 1
    \node[oplabel] (op1) at (7.2, -1.8) {$\hat{\mathcal{P}}_{\Delta t}$};

    % =========================================================================
    % Nivel 1: Muestreo Simultáneo (y = -3.6)
    % =========================================================================
    \node[block] (b1) at (0, -3.6) {
        \textbf{Muestreo simultáneo en FPGA}\\
        \footnotesize Cuantización determinista ($\tau_{\text{skew}} \equiv 0$)
    };
    \node[state] (s1) at (7.2, -3.6) {
        $\mathbf{p}[n] = \begin{pmatrix} p_1[n\Delta t] \pm \Delta p \\ p_2[n\Delta t] \pm \Delta p \end{pmatrix}$
    };

    % Operador 2
    \node[oplabel] (op2) at (7.2, -5.4) {$\hat{\mathcal{C}}_M$};

    % =========================================================================
    % Nivel 2: Diezmado (y = -7.2)
    % =========================================================================
    \node[block] (b2) at (0, -7.2) {
        \textbf{Diezmado digital promediado}\\
        \footnotesize Reducción estocástica de varianza
    };
    \node[state] (s2) at (7.2, -7.2) {
        $\tilde{\mathbf{p}}[k] = \begin{pmatrix} \tilde{p}_1[k\Delta t'] \pm \frac{\Delta p}{\sqrt{M}} \\ \tilde{p}_2[k\Delta t'] \pm \frac{\Delta p}{\sqrt{M}} \end{pmatrix}$
    };

    % Operador 3
    \node[oplabel] (op3) at (7.2, -9.0) {$\hat{\mathcal{W}}_{nc}$};

    % =========================================================================
    % Nivel 3: Enventanado Causal (y = -10.8)
    % =========================================================================
    \node[block] (b3) at (0, -10.8) {
        \textbf{Enventanado temporal causal}\\
        \footnotesize Truncamiento a soporte compacto $\mathbb{R}^N$
    };
    \node[state] (s3) at (7.2, -10.8) {
        $\vec{P} = \begin{pmatrix} \vec{p}_1 \pm \Delta\vec{p}' \\ \vec{p}_2 \pm \Delta\vec{p}' \end{pmatrix} \in \mathbb{R}^N \times \mathbb{R}^N$
    };

    % =========================================================================
    % CONEXIONES VERTICALES INTERNAS (Etapa 1)
    % =========================================================================
    \draw[arrow] (b0) -- (b1);
    \draw[arrow] (b1) -- (b2);
    \draw[arrow] (b2) -- (b3);

    \draw[arrow] (s0) -- (op1);
    \draw[arrow] (op1) -- (s1);

    \draw[arrow] (s1) -- (op2);
    \draw[arrow] (op2) -- (s2);

    \draw[arrow] (s2) -- (op3);
    \draw[arrow] (op3) -- (s3);

    % =========================================================================
    % BIFURCACIÓN ELEGANTE (Convergencia suave y tronco común con etiqueta)
    % =========================================================================
    % 1. Curvas suaves de unión hacia el centro (x = 3.6, y = -12.4)
    \coordinate (merge) at (3.6, -12.4);
    \draw[line] (b3.south) to[out=-90, in=180] (merge);
    \draw[line] (s3.south) to[out=-90, in=0] (merge);

    % 2. Tronco común vertical con etiqueta física del vector que se bifurca
    \coordinate (split) at (3.6, -13.5);
    \draw[line] (merge) -- node[right=5pt, font=\small, align=left] {
        \textbf{Vector de estado $\vec{P}$}\\
        \footnotesize (Acondicionado)
    } (split);

    % 3. Cajas de destino de las ramas (y = -14.8)
    \node[branchbox] (ramaA) at (0.0, -14.8) {
        \textbf{Rama A: Canal temporal}\\
        \footnotesize Banda estrecha en $L^2$ (Sección~1.3)
    };
    \node[branchbox] (ramaB) at (7.2, -14.8) {
        \textbf{Rama B: Espacio recíproco}\\
        \footnotesize Descomposición en $\ell^2$ (Sección~1.4)
    };

    % 4. Flechas ortogonales limpias hacia cada rama
    \draw[arrow] (split) -| (ramaA.north);
    \draw[arrow] (split) -| (ramaB.north);

\end{tikzpicture}